\chapter[1912 --- Carrel]{1912: Alexis Carrel}
\label{ch:1912_alexiscarrel}

\section*{Main Contribution}

\textit{Developed techniques for suturing blood vessels and transplanting organs, laying the groundwork for modern vascular and transplant surgery.}

\section{Official Citation}

for his work on vascular suture and the transplantation of blood vessels and organs

\section{Background and Historical Context}

Alexis Carrel was born on June 28, 1873, in Sainte-Foy-l\`es-Lyon, a suburb of Lyon, France. The son of a prosperous manufacturing family, Carrel studied medicine at the University of Lyon, receiving his medical degree in 1900. He quickly became interested in surgery, particularly the techniques of suturing blood vessels, and devoted himself to developing methods for the repair of damaged tissues. After a period of research at the University of Chicago and at the Rockefeller Institute for Medical Research in New York, Carrel joined the Rockefeller Institute permanently in 1906, where he conducted his primary work.

The state of surgery in the early twentieth century was advancing rapidly but was still limited by fundamental technical challenges. The development of aseptic technique and anesthesia had made surgery safer, and surgeons were increasingly ambitious in the scope of their operations. But the repair of damaged blood vessels---the essential highways of the circulatory system---remained a formidable challenge. When a major blood vessel was severed or damaged, the surgeon had few options: the vessel could be ligated (tied off), but this cut off blood supply to the tissues beyond the injury, often leading to gangrene and amputation. The ability to repair damaged blood vessels by suturing them back together---vascular anastomosis---was the holy grail of surgical technique.

Prior to Carrel's work, several surgeons had attempted to suture blood vessels, with limited success. The Italian surgeon Angelo Saporiti had performed experimental vessel repair in animals in the 1890s, and the French surgeon Jean-Louis Chaput had developed techniques for vascular anastomosis. But these methods were technically demanding, unreliable, and often resulted in thrombosis (clotting) or leakage at the repair site.

The transplantation of organs presented an even greater challenge. The kidney, heart, liver, and other organs could not be removed from one individual and transplanted into another without first being able to reconnect the blood vessels that supplied them. Without reliable techniques for vascular anastomosis, organ transplantation was a theoretical impossibility.

\section{The Discovery}

Carrel's Nobel Prize-winning work addressed these challenges through a combination of surgical innovation, experimental ingenuity, and careful study of the biological processes of wound healing.

\subsection{Vascular Suture}

Carrel's primary contribution was the development of a reliable technique for suturing blood vessels---the carotid-carotid anastomosis technique that became the standard method for vascular repair. His technique involved several key innovations:

\textbf{Fine suture material.} Carrel used the finest available silk sutures, cut into short lengths, and threaded through delicate, atraumatic needles. The use of fine suture material minimized tissue trauma and reduced the risk of thrombosis at the suture line.

\textbf{Eversion of vessel edges.} Carrel's technique involved turning the cut edges of the vessel outward (everting) before suturing, so that the inner lining of the vessel (the intima) was pressed against the intima of the opposing vessel wall. This eversion technique ensured that the smooth, thrombosis-resistant intimal surfaces were in contact at the suture line, reducing the risk of clot formation.

\textbf{Three-suture technique.} Carrel developed a method in which three evenly spaced traction sutures were placed through the full thickness of the vessel wall at the edges of the cut. These sutures were then used to rotate the vessel during suturing, ensuring an even and airtight closure. This technique simplified the operation and improved the consistency of results.

\textbf{Aseptic technique.} Carrel insisted on strict aseptic technique throughout the operation, recognizing that infection at the suture line would lead to failure.

Through these innovations, Carrel achieved consistently successful vascular anastomosis in experimental animals. He demonstrated that arteries and veins could be sutured together, that repaired vessels healed and remained patent (open) for months or years, and that blood flow through repaired vessels was restored to normal. His publications on vascular suture, beginning in 1902, were widely read and rapidly adopted by surgeons worldwide.

\subsection{Organ Transplantation}

Carrel's techniques for vascular anastomosis opened the door to organ transplantation---the removal of an organ from one individual and its transplantation into another. In a series of landmark experiments, Carrel demonstrated that kidneys could be transplanted from one dog to another, with restoration of function after the vascular connections were re-established. He showed that the transplanted organ survived and produced urine, confirming that vascular anastomosis made organ transplantation technically feasible.

Carrel also conducted the first experimental transplantation of the heart, transplanting the heart of a small dog into the neck of a larger dog and demonstrating that the transplanted heart continued to beat and circulate blood. While cardiac transplantation would not become clinically feasible for another half-century, Carrel's experiments demonstrated the principle and laid the groundwork for the eventual development of the field.

\subsection{Culture of Tissues}

Beyond his surgical innovations, Carrel made important contributions to the study of tissue culture. Working with the American embryologist and tissue culture pioneer Montrose Thomas Burrows, Carrel developed techniques for the cultivation of cells and tissues outside the body. He demonstrated that cells from chicken hearts could be maintained in culture indefinitely, provided that the culture medium was changed regularly and that the cells were periodically transferred to fresh culture vessels. This was one of the first demonstrations of continuous cell culture---the maintenance of living cells in the laboratory---and it had far-reaching implications for cell biology, cancer research, and the development of vaccines and other biological products.

Carrel's tissue culture work also led him to study the biology of aging. He proposed that cells had an inherent capacity for unlimited division if provided with optimal culture conditions---a claim that was controversial at the time and was later shown to be incorrect (normal somatic cells have a finite replicative capacity, known as the Hayflick limit). Nevertheless, Carrel's tissue culture techniques were adopted and refined by subsequent researchers and became essential tools of modern cell biology.

\section{Impact on Modern Medicine}

Carrel's contributions to vascular surgery and organ transplantation have had a transformative impact on medicine. The techniques of vascular anastomosis that he developed are still used, in refined forms, in vascular surgery, cardiac surgery, and transplant surgery today. The microsurgical techniques used in modern vascular and reconstructive surgery are direct descendants of Carrel's pioneering work.

Organ transplantation, which Carrel made technically feasible, is now a standard medical procedure that saves tens of thousands of lives each year. Kidney, liver, heart, lung, and pancreas transplants are performed in hospitals around the world, and the development of immunosuppressive drugs has made it possible to transplant organs between unrelated donors with acceptable rates of graft survival. The first successful human kidney transplant was performed by Joseph Murray in 1954, and the first successful human heart transplant was performed by Christiaan Barnard in 1967---achievements that were made possible by the surgical techniques that Carrel had developed decades earlier.

Carrel's tissue culture techniques have had an equally significant impact. The maintenance and manipulation of cells in culture is now a fundamental tool of biomedical research, used in the study of cancer, infectious disease, drug development, and regenerative medicine. The HeLa cell line, derived from cervical cancer cells taken from Henrietta Lacks in 1951, is a direct descendant of the continuous cell culture techniques that Carrel developed. Modern stem cell research, tissue engineering, and the development of organoids (miniature organ-like structures grown from stem cells) all build on the tissue culture foundations that Carrel established.

\section{Legacy}

Alexis Carrel received the Nobel Prize in Physiology or Medicine in 1912 ``for his work on vascular suture and the transplantation of blood vessels and organs.'' He was 39 years old at the time of the award, one of the younger recipients of the prize.

Carrel continued his research at the Rockefeller Institute until 1939, when he returned to France during the German occupation. He worked in Paris under the Vichy government, a controversial association that has clouded his legacy. He died on November 5, 1944, at the age of 71.

Despite the controversies of his later years, Carrel's contributions to surgery and transplantation are beyond dispute. His techniques for vascular anastomosis and organ transplantation saved countless lives and opened new frontiers in surgical practice. His tissue culture methods provided tools that remain essential for biomedical research. And his vision of organ transplantation---which seemed impossibly futuristic in the early twentieth century---is now a routine reality of modern medicine.


\section{Modern Developments}

Carrel's vascular surgery techniques evolved into modern transplant surgery. Hand transplantation (first performed 1998) and face transplantation (2005) build on his vascular anastomosis methods. Organ preservation solutions (University of Wisconsin solution) extend transplant organ viability to 36 hours. Machine perfusion systems now keep donor hearts beating outside the body for hours, improving transplant outcomes. CRISPR gene editing is being used to create xenotransplantation-compatible pigs for organ transplantation to humans.


\section{Bibliography}

\begin{thebibliography}{9}
\bibitem{nobel-1912}
Nobel Prize Outreach, ``The Nobel Prize in Physiology or Medicine 1912,''\emph{NobelPrize.org}.\\
\url{https://www.nobelprize.org/prizes/medicine/1912/summary/}

\bibitem{lecture-1912}
Alexis Carrel, ``Nobel Lecture,''\emph{NobelPrize.org}. Winner-authored primary source; downloadable from the official Nobel Prize archive.\\
\url{https://www.nobelprize.org/prizes/medicine/1912/carrel/lecture/}
\end{thebibliography}
